\documentclass[10pt]{article}
\usepackage[margin=1in]{geometry}
\usepackage{amsmath, amssymb}
\usepackage{enumitem}

\title{Lesson Plan: Properties of Logarithms}
\author{}
\date{}

\begin{document}

\maketitle

\section*{Objective}
Students will understand and apply the properties of logarithms to expand, condense, and evaluate logarithmic expressions.

\section*{Standards}
\begin{itemize}
    \item Rewrite logarithmic expressions using properties of logarithms.
    \item Solve problems involving exponential and logarithmic relationships.
\end{itemize}

\section*{Prerequisites}
\begin{itemize}
    \item Understanding of exponents and exponent rules
    \item Definition of logarithms as inverses of exponential functions
\end{itemize}

\section*{Materials}
\begin{itemize}
    \item Whiteboard / projector
    \item Student worksheet
    \item Graphing calculator (optional)
\end{itemize}

\section*{Lesson Duration}
45--60 minutes

\section*{Lesson Structure}

\subsection*{1. Warm-Up (5--10 minutes)}

Rewrite using exponent rules:
\begin{enumerate}[label=\arabic*.]
    \item $a^m \cdot a^n =$
    \item $\frac{a^m}{a^n} =$
    \item $(a^m)^n =$
\end{enumerate}

Then ask:
\begin{quote}
How might these rules translate to logarithms?
\end{quote}

\subsection*{2. Motivation (5--10 minutes)}

Start from the definition:
\[
\log_b(x) = y \iff b^y = x
\]

Example:
\[
\log_2(8) = 3 \quad \text{because} \quad 2^3 = 8
\]

Pose the question:
\begin{quote}
What is $\log_b(xy)$ in terms of $\log_b(x)$ and $\log_b(y)$?
\end{quote}

\subsection*{3. Direct Instruction (15 minutes)}

\textbf{Product Rule:}
\[
\log_b(xy) = \log_b(x) + \log_b(y)
\]

\textbf{Quotient Rule:}
\[
\log_b\left(\frac{x}{y}\right) = \log_b(x) - \log_b(y)
\]

\textbf{Power Rule:}
\[
\log_b(x^k) = k \log_b(x)
\]

\textbf{Important Conditions:}
\begin{itemize}
    \item $x > 0$, $y > 0$
    \item $b > 0$, $b \neq 1$
\end{itemize}

\textbf{Example 1 (Expansion):}
\[
\log_2(8x^3) = \log_2(8) + \log_2(x^3) = 3 + 3\log_2(x)
\]

\textbf{Example 2 (Condensing):}
\[
2\log(x) - \log(y) = \log(x^2) - \log(y) = \log\left(\frac{x^2}{y}\right)
\]

\subsection*{4. Guided Practice (10--15 minutes)}

\begin{enumerate}[label=\arabic*.]
    \item Expand: $\log_3(9x^2y)$
    \item Expand: $\log\left(\frac{a^3}{b\sqrt{c}}\right)$
    \item Condense: $\log(x) + \log(y) - \log(z)$
    \item Condense: $3\log(a) + \frac{1}{2}\log(b)$
\end{enumerate}

Encourage students to justify each step using a named property.

\subsection*{5. Independent Practice (10 minutes)}

\begin{enumerate}[label=\arabic*.]
    \item Expand completely: $\log_5\left(\frac{25x^4}{y^2}\right)$
    \item Condense into a single logarithm: 
    \[
    2\log(x) - 3\log(y) + \log(z)
    \]
    \item Evaluate: $\log_2(32) + \log_2(4)$
    \item Challenge: Solve for $x$:
    \[
    \log(x) + \log(x - 3) = 1
    \]
\end{enumerate}

\subsection*{6. Closure (5 minutes)}

Ask:
\begin{itemize}
    \item How are logarithm properties related to exponent rules?
    \item Why must arguments of logarithms be positive?
\end{itemize}

Summarize:
\begin{quote}
Logarithm properties transform multiplication into addition, division into subtraction, and exponents into coefficients.
\end{quote}

\section*{Assessment}

\begin{itemize}
    \item Exit Ticket:
    \[
    \text{Condense: } \log(x) + 2\log(y) - \log(z)
    \]
\end{itemize}

\section*{Differentiation}

\textbf{Support:}
\begin{itemize}
    \item Provide a reference sheet of exponent rules alongside log rules
    \item Use numerical examples before variables
\end{itemize}

\textbf{Extension:}
\begin{itemize}
    \item Introduce change of base formula
    \item Solve logarithmic equations with multiple steps
\end{itemize}

\end{document}